%% ========================================================================
%% Chapter 9: Pemrograman Bash
%% ========================================================================

\chapter{Pemrograman Bash}
\label{ch:bash-programming}

\chapterhero{purple}{Ubah rangkaian perintah manual menjadi skrip yang rapi, dapat dipakai ulang, dan lebih aman melalui fungsi, validasi, serta debugging Bash.}{\faFileCode\quad script}{\faPuzzlePiece\quad function}{\faBug\quad debugging}

Di Bab~6 kamu belajar cara bekerja lebih cepat di terminal: alias,
fungsi shell sederhana, variabel lingkungan, dan ekspansi. Semua itu
adalah bekal untuk bab ini. Pemrograman Bash adalah langkah berikutnya:
menyimpan serangkaian perintah ke dalam file, menambahkan logika
pengambilan keputusan, dan menjalankannya otomatis tanpa campur tangan
manusia.

Bayangkan kamu setiap pagi membuka terminal dan menjalankan serangkaian
perintah yang sama: cek disk, cek memori, arsipkan log kemarin, kirim
laporan. Membosankan dan rawan lupa. Sebuah skrip Bash mengotomatiskan
semua itu dengan satu perintah atau bahkan tanpa kamu harus hadir sama
sekali. Skrip Bash adalah ``resep masakan'' yang bisa dijalankan
kapan saja, oleh siapa saja yang punya akses yang sama, dan menghasilkan
hasil yang konsisten setiap kali.

Bab ini membangun langsung di atas fondasi Bash dari Bab~6.
Di bab ini kamu akan menulis skrip dari nol, mengelola variabel dan
parameter, menyusun logika percabangan dan perulangan, membangun fungsi
yang bisa dipakai ulang, dan akhirnya men-debug skrip yang bermasalah.
Semua teknik ini menjadi fondasi untuk Bab~8 di mana kamu akan membuat
skrip pemantauan memori sistem yang sesungguhnya.

\textbf{Yang Akan Kamu Pelajari}

Setelah menyelesaikan bab ini, kamu akan mampu:

\begin{enumerate}
    \item membuat skrip Bash yang bisa dieksekusi dengan shebang dan izin
          yang benar, serta memahami peran exit code;
    \item menggunakan variabel, parameter posisional \texttt{\$1},
          \texttt{\$\#}, \texttt{\$@}, dan aritmetika Bash;
    \item menulis struktur kontrol \cmd{if}/\cmd{elif}/\cmd{else},
          \cmd{for}, \cmd{while}, dan \cmd{case} untuk mengotomatiskan
          logika pengambilan keputusan;
    \item mendefinisikan fungsi shell dengan variabel lokal dan memisahkan
          fungsi ke file library yang bisa di-\cmd{source};
    \item men-debug skrip menggunakan \cmd{set -x}, \cmd{set -e},
          \cmd{set -u}, dan \cmd{trap} untuk cleanup otomatis.
\end{enumerate}

\begin{notebox}
Seluruh praktek pada bab ini dijalankan pada direktori
\dir{~/lab-os/chapter07-scripting/}.
Buat direktori ini sebelum memulai:
\begin{lstlisting}[language=bash]
mkdir -p ~/lab-os/chapter07-scripting/{skrip,logs,lib}
cd ~/lab-os/chapter07-scripting
\end{lstlisting}
\end{notebox}

%% ------------------------------------------------------------------------
\section{Dasar Skrip Bash}
\label{sec:dasar-skrip-bash}

Skrip Bash adalah file teks biasa yang berisi perintah-perintah yang
biasa kamu ketik di terminal, disusun berurutan. Saat dijalankan, Bash
membaca file tersebut baris per baris dan mengeksekusinya satu per satu.

Baris pertama skrip hampir selalu berupa \textit{shebang}: \texttt{\#!/bin/bash}.
Karakter \texttt{\#!} memberi tahu kernel Linux bahwa file ini adalah
skrip dan \texttt{/bin/bash} adalah interpreter yang harus menjalankannya.
Tanpa shebang, sistem tidak tahu program apa yang harus memproses file.
Kamu masih bisa menjalankannya dengan \texttt{bash skrip.sh}, tetapi cara
itu melewati shebang secara eksplisit. Alternatif portabel yang umum
dipakai untuk skrip yang perlu berjalan di berbagai sistem adalah
\texttt{\#!/usr/bin/env bash}, yang menyuruh sistem mencari Bash di
\var{PATH}.

Setelah menulis shebang, skrip perlu izin eksekusi sebelum bisa
dijalankan langsung. Berikan dengan \texttt{chmod +x nama-skrip.sh},
lalu jalankan dengan \texttt{./nama-skrip.sh}. Awalan \texttt{./} penting:
ia memberitahu shell untuk mencari file di direktori saat ini, bukan di
\var{PATH}.

Setiap perintah di Bash menghasilkan \textit{exit code}, yaitu angka yang
menandakan apakah perintah berhasil atau gagal. Konvensi universal: 0
berarti sukses, angka 1 sampai 255 berarti ada kesalahan. Variabel
\texttt{\$?} menyimpan exit code perintah terakhir. Kamu bisa memeriksa
exit code untuk membuat skrip yang bereaksi terhadap kegagalan.

Komentar di Bash dimulai dengan tanda \texttt{\#}. Gunakan komentar untuk
menjelaskan tujuan skrip, parameter yang dibutuhkan, dan logika yang tidak
langsung jelas. Skrip tanpa komentar sulit dipahami saat kamu kembali
melihatnya enam bulan kemudian.

\subsection*{Praktek 7.1: Buat dan Jalankan Skrip Pertama}

Tujuan praktek ini adalah membuat skrip Bash lengkap dengan shebang,
memberikan izin eksekusi, dan memahami exit code.

\begin{enumerate}
    \item Buat skrip pertama dengan pola \textit{Hello World} yang juga menampilkan informasi sistem dasar:
\begin{lstlisting}[language=bash, caption={Membuat skrip hello-world.sh}]
nano ~/lab-os/chapter07-scripting/skrip/hello-world.sh
# Tulis isi berkas berikut
#!/bin/bash
# hello-world.sh: Menampilkan Hello World dan info sistem dasar
# Penggunaan: ./hello-world.sh [nama-pengguna]

NAMA=${1:-"Tamu"}

echo "==============================="
echo "  Hello, World!"
echo "  Halo, $NAMA!"
echo "==============================="
echo "Hostname : $(hostname)"
echo "User     : $(whoami)"
echo "Tanggal  : $(date '+%A, %d %B %Y')"
echo "Jam      : $(date '+%H:%M:%S')"
echo "Uptime   : $(uptime -p)"
echo "==============================="
\end{lstlisting}

    \item Berikan izin eksekusi dan jalankan:
\begin{lstlisting}[language=bash, caption={Memberikan izin dan menjalankan skrip}]
chmod +x ~/lab-os/chapter07-scripting/skrip/hello-world.sh

cd ~/lab-os/chapter07-scripting/skrip
./hello-world.sh
./hello-world.sh "Budi"
\end{lstlisting}
    \textit{Amati:} tanpa argumen, skrip menggunakan ``Tamu'' sebagai
    nama default. Dengan argumen, skrip menggunakan nama yang diberikan.
    Ini karena konstruk \texttt{\$\{1:-"Tamu"\}} membaca \texttt{\$1};
    jika kosong, nilai default digunakan.

    \item Periksa exit code skrip:
\begin{lstlisting}[language=bash, caption={Memeriksa exit code}]
./hello-world.sh "Ani"
echo "Exit code: $?"

./hello-world.sh "Budi" > /dev/null
echo "Exit code setelah redirect: $?"
\end{lstlisting}
    \textit{Amati:} exit code 0 menandakan skrip berjalan tanpa kesalahan.
    Redirect ke \file{/dev/null} membuang output ke tempat kosong, berguna
    saat kamu hanya peduli dengan exit code bukan outputnya.

    \item Tambahkan exit code eksplisit ke skrip dengan kasus error:
\begin{lstlisting}[language=bash, caption={Menambahkan penanganan error sederhana}]
nano ~/lab-os/chapter07-scripting/skrip/cek-file.sh
# Tulis isi berkas berikut
#!/bin/bash
# cek-file.sh: Memeriksa apakah file ada dan bisa dibaca

if [ $# -ne 1 ]; then
    echo "Penggunaan: $0 <path-file>"
    exit 1
fi

if [ -f "$1" ] && [ -r "$1" ]; then
    echo "File ditemukan: $1"
    echo "Ukuran: $(wc -c < "$1") byte"
    exit 0
else
    echo "File tidak ditemukan atau tidak bisa dibaca: $1"
    exit 2
fi

chmod +x ~/lab-os/chapter07-scripting/skrip/cek-file.sh
./cek-file.sh /etc/hostname
echo "Exit code: $?"

./cek-file.sh /file-tidak-ada.txt
echo "Exit code: $?"
\end{lstlisting}
\end{enumerate}

\begin{challengebox}
Buat skrip \cmd{sapa-formal.sh} yang menerima nama pengguna sebagai
argumen. Skrip harus menampilkan salam yang berbeda berdasarkan jam saat
ini: ``Selamat pagi'' untuk jam 05.00 sampai 11.59, ``Selamat siang''
untuk jam 12.00 sampai 14.59, ``Selamat sore'' untuk jam 15.00 sampai
17.59, dan ``Selamat malam'' untuk jam lainnya. Gunakan command
substitution \texttt{\$(date +\%H)} untuk mendapatkan jam saat ini
sebagai angka. Jika tidak ada argumen, skrip menampilkan pesan penggunaan
dan keluar dengan exit code 1.
\end{challengebox}

%% ------------------------------------------------------------------------
\section{Variabel dan Parameter Posisional}
\label{sec:variabel-parameter-posisional}

Variabel di Bash bekerja berbeda dari bahasa pemrograman yang mengenal
tipe data. Semua variabel Bash adalah string. Saat kamu
menulis \texttt{ANGKA=42}, Bash menyimpan teks ``42'', bukan bilangan
bulat. Operasi aritmetika memerlukan sintaks khusus.

Aturan penting saat membuat variabel: tidak boleh ada spasi di sekitar
tanda sama dengan. \texttt{NAMA="Budi"} benar. \texttt{NAMA = "Budi"}
salah karena Bash mengira \texttt{NAMA} adalah nama perintah yang dipanggil
dengan argumen \texttt{=} dan \texttt{"Budi"}.

Untuk membaca nilai variabel, gunakan \texttt{\$NAMA} atau
\texttt{\$\{NAMA\}}. Bentuk kurung kurawal berguna saat variabel disambung
langsung dengan teks, misalnya \texttt{\$\{NAMA\}\_backup.txt} agar Bash
tahu batas nama variabelnya.

Parameter posisional adalah variabel otomatis yang mewakili argumen
yang diberikan saat skrip dijalankan. \texttt{\$0} adalah nama skrip
itu sendiri. \texttt{\$1} adalah argumen pertama, \texttt{\$2} argumen
kedua, dan seterusnya. \texttt{\$\#} menyimpan jumlah argumen.
\texttt{\$@} mewakili semua argumen sebagai daftar yang bisa diiterasi,
sedangkan \texttt{\$*} menggabungkan semua argumen menjadi satu string.
\texttt{\$?} adalah exit code perintah terakhir. \texttt{\$\$} adalah
PID (Process ID) skrip yang sedang berjalan.

Aritmetika Bash menggunakan sintaks \texttt{\$(( ekspresi ))}. Di
dalamnya kamu bisa menggunakan \texttt{+}, \texttt{-}, \texttt{*},
\texttt{/} (pembagian bulat), dan \texttt{\%} (modulo). Untuk desimal,
gunakan perintah \cmd{bc}.

\subsection*{Praktek 7.2: Gunakan Parameter Posisional}

Tujuan praktek ini adalah membuat skrip yang memproses argumen dari
command line dan menggunakan aritmetika.

\begin{enumerate}
    \item Buat skrip yang menampilkan semua parameter posisional:
\begin{lstlisting}[language=bash, caption={Membuat skrip cek-parameter.sh}]
nano ~/lab-os/chapter07-scripting/skrip/cek-parameter.sh
# Tulis isi berkas berikut
#!/bin/bash
# cek-parameter.sh: Menampilkan informasi parameter posisional

echo "Nama skrip  : $0"
echo "Arg pertama : $1"
echo "Arg kedua   : $2"
echo "Jumlah arg  : $#"
echo "Semua arg   : $@"

echo ""
echo "Iterasi semua argumen:"
NOMOR=1
for ARG in "$@"; do
    echo "  Arg $NOMOR: $ARG"
    NOMOR=$((NOMOR + 1))
done

chmod +x ~/lab-os/chapter07-scripting/skrip/cek-parameter.sh
./cek-parameter.sh apel mangga "jeruk nipis" pepaya
\end{lstlisting}
    \textit{Amati:} ``jeruk nipis'' diperlakukan sebagai satu argumen
    karena dikutip. Tanpa kutipan, Bash memecahnya menjadi dua argumen
    terpisah.

    \item Buat skrip yang menghitung statistik sederhana dari daftar angka:
\begin{lstlisting}[language=bash, caption={Membuat skrip statistik-angka.sh}]
nano ~/lab-os/chapter07-scripting/skrip/statistik-angka.sh
# Tulis isi berkas berikut
#!/bin/bash
# statistik-angka.sh: Hitung jumlah dan rata-rata dari argumen angka

if [ $# -eq 0 ]; then
    echo "Penggunaan: $0 <angka1> <angka2> ..."
    echo "Contoh: $0 10 20 30 40 50"
    exit 1
fi

JUMLAH=0
BANYAK=$#

for ANGKA in "$@"; do
    JUMLAH=$((JUMLAH + ANGKA))
done

RATA=$((JUMLAH / BANYAK))

echo "Banyak data : $BANYAK"
echo "Total       : $JUMLAH"
echo "Rata-rata   : $RATA (pembulatan ke bawah)"

chmod +x ~/lab-os/chapter07-scripting/skrip/statistik-angka.sh
./statistik-angka.sh 10 20 30 40 50
./statistik-angka.sh 7 14 21
\end{lstlisting}

    \item Uji penanganan error saat tidak ada argumen:
\begin{lstlisting}[language=bash, caption={Menguji penanganan argumen kosong}]
./statistik-angka.sh
echo "Exit code: $?"
\end{lstlisting}
    \textit{Amati:} skrip menampilkan pesan penggunaan dan keluar dengan
    exit code 1, bukan crash dengan pesan error yang membingungkan.
\end{enumerate}

\begin{challengebox}
Buat skrip \cmd{kalkulator.sh} yang menerima tiga argumen: angka pertama,
operator (\texttt{+}, \texttt{-}, \texttt{x}, atau \texttt{/}), dan angka
kedua. Gunakan \cmd{case} untuk memilih operasi yang sesuai. Tambahkan
validasi: jika jumlah argumen bukan 3, tampilkan pesan penggunaan dan exit
code 1. Jika operator tidak dikenal, tampilkan pesan error dan exit code 2.
Untuk pembagian, tangani kasus pembagi nol secara khusus dengan exit code
3. Contoh pemanggilan: \texttt{./kalkulator.sh 20 + 5} menghasilkan
\texttt{Hasil: 25}, dan \texttt{./kalkulator.sh 10 / 0} menghasilkan
pesan ``Error: pembagi tidak boleh nol''.
\end{challengebox}

%% ------------------------------------------------------------------------
\section{Struktur Kontrol}
\label{sec:struktur-kontrol-bash}

Struktur kontrol membuat skrip bisa membuat keputusan dan mengulang
tindakan. Tanpa struktur kontrol, skrip hanya menjalankan perintah secara
berurutan dari atas ke bawah. Dengan struktur kontrol, skrip bisa
bereaksi berbeda terhadap kondisi yang berbeda.

Perintah \cmd{if} di Bash menggunakan perintah \cmd{test} atau tanda
kurung siku \texttt{[ kondisi ]} untuk mengevaluasi kondisi. Tanda
kurung siku adalah alias untuk perintah \cmd{test}. Bash juga
mengenal \texttt{[[ kondisi ]]} yang merupakan versi lebih modern dengan
lebih banyak fitur, termasuk pencocokan pola dengan \texttt{=\textasciitilde}.

Operator perbandingan untuk angka: \texttt{-eq} (sama), \texttt{-ne}
(tidak sama), \texttt{-gt} (lebih besar), \texttt{-lt} (lebih kecil),
\texttt{-ge} (lebih besar atau sama), \texttt{-le} (lebih kecil atau sama).
Untuk string: \texttt{=} atau \texttt{==} (sama), \texttt{!=} (tidak sama),
\texttt{-z} (string kosong), \texttt{-n} (string tidak kosong).
Untuk file: \texttt{-f} (file biasa ada), \texttt{-d} (direktori ada),
\texttt{-r} (bisa dibaca), \texttt{-w} (bisa ditulis), \texttt{-x}
(bisa dieksekusi).

Perulangan \cmd{for} di Bash bisa mengiterasi daftar nilai atau rentang
angka. Bentuk \texttt{for VAR in daftar; do ... done} mengiterasi setiap
item dalam daftar. Bentuk C-style \texttt{for ((i=1; i<=5; i++)); do ...
done} mirip dengan bahasa pemrograman lain. Perulangan \cmd{while}
terus berjalan selama kondisi bernilai benar. Perulangan \cmd{until}
adalah kebalikannya: berjalan sampai kondisi bernilai benar.

Perintah \cmd{case} adalah cara elegan untuk menangani banyak kondisi
berbeda berdasarkan nilai satu variabel. Setiap blok pola diakhiri
\texttt{;;} untuk menghentikan eksekusi. Pola \texttt{*)} adalah tangkap
semua yang cocok dengan kondisi apa pun.

\subsection*{Praktek 7.3: Tulis Skrip dengan Logika Percabangan dan Perulangan}

Tujuan praktek ini adalah menggabungkan \cmd{if}, \cmd{for}, \cmd{while},
dan \cmd{case} dalam skrip yang berguna secara nyata.

\begin{enumerate}
    \item Buat skrip yang memeriksa kondisi setiap filesystem:
\begin{longlisting}[language=bash, caption={Membuat skrip cek-disk.sh}]
nano ~/lab-os/chapter07-scripting/skrip/cek-disk.sh
# Tulis isi berkas berikut
#!/bin/bash
# cek-disk.sh: Periksa penggunaan disk setiap filesystem
# Penggunaan: ./cek-disk.sh [batas-persen]

BATAS=${1:-80}

echo "=== CEK DISK ==="
echo "Batas peringatan: ${BATAS}%"
echo ""

PERINGATAN=0

while IFS= read -r BARIS; do
    PERSEN=$(echo "$BARIS" | awk '{print $5}' | tr -d '%')
    MOUNT=$(echo "$BARIS" | awk '{print $6}')
    PERANGKAT=$(echo "$BARIS" | awk '{print $1}')

    if [ "$PERSEN" -ge "$BATAS" ]; then
        echo "[PENUH]  $MOUNT ($PERSEN%) - $PERANGKAT"
        PERINGATAN=$((PERINGATAN + 1))
    elif [ "$PERSEN" -ge $((BATAS / 2)) ]; then
        echo "[SEDANG] $MOUNT ($PERSEN%) - $PERANGKAT"
    else
        echo "[AMAN]   $MOUNT ($PERSEN%) - $PERANGKAT"
    fi
done < <(df -h | tail -n +2)

echo ""
echo "Total peringatan: $PERINGATAN filesystem"

chmod +x ~/lab-os/chapter07-scripting/skrip/cek-disk.sh
./cek-disk.sh
./cek-disk.sh 50
\end{longlisting}
    \textit{Amati:} konstruk \texttt{while IFS= read -r BARIS} membaca
    output perintah baris per baris. Ini pola yang sangat umum di skrip
    Bash untuk memproses output perintah lain.

    \item Buat skrip dengan \cmd{case} untuk menu monitoring:
\begin{longlisting}[language=bash, caption={Membuat skrip menu-monitor.sh}]
nano ~/lab-os/chapter07-scripting/skrip/menu-monitor.sh
# Tulis isi berkas berikut
#!/bin/bash
# menu-monitor.sh: Menu interaktif pemantauan sistem

while true; do
    echo ""
    echo "=== MENU MONITOR SISTEM ==="
    echo "1) Status disk"
    echo "2) Status memori"
    echo "3) Proses teratas"
    echo "4) Info jaringan"
    echo "5) Keluar"
    echo -n "Pilih [1-5]: "
    read -r PILIHAN

    case $PILIHAN in
        1)
            echo "--- Status Disk ---"
            df -h
            ;;
        2)
            echo "--- Status Memori ---"
            free -h
            ;;
        3)
            echo "--- 5 Proses Teratas (CPU) ---"
            ps aux --sort=-%cpu | head -6
            ;;
        4)
            echo "--- Info Jaringan ---"
            ip -brief addr show 2>/dev/null || ifconfig 2>/dev/null | head -20
            ;;
        5)
            echo "Sampai jumpa!"
            exit 0
            ;;
        *)
            echo "Pilihan tidak valid. Masukkan angka 1-5."
            ;;
    esac
done

chmod +x ~/lab-os/chapter07-scripting/skrip/menu-monitor.sh
\end{longlisting}
    \textit{Amati:} jalankan skrip dengan \texttt{./menu-monitor.sh} dan
    coba setiap pilihan. Tekan \keystroke{Ctrl+C} atau pilih 5 untuk
    keluar. Perhatikan bagaimana \cmd{while true} membuat menu terus
    muncul sampai pengguna memilih keluar.
\end{enumerate}

\begin{challengebox}
Buat skrip \cmd{cek-etc.sh} yang memeriksa setiap direktori langsung di
dalam \dir{/etc} dan melaporkan apakah direktori tersebut kosong atau
berisi file. Gunakan perintah \cmd{find} dari Bab~4 untuk memeriksa isi
direktori. Tampilkan dua kolom: nama direktori dan statusnya (KOSONG atau
BERISI diikuti jumlah file). Gunakan \cmd{for} untuk mengiterasi direktori
dan \cmd{if} untuk memeriksa kondisi. Simpan laporan ke
\file{~/lab-os/chapter07-scripting/logs/laporan-etc.txt} menggunakan
redirection yang dipelajari di Bab~3, sekaligus tampilkan ke terminal
dengan \cmd{tee}.
\end{challengebox}

%% ------------------------------------------------------------------------
\section{Fungsi Shell}
\label{sec:fungsi-shell-skrip}

Fungsi adalah blok kode yang diberi nama dan bisa dipanggil berulang kali.
Bayangkan kamu menulis skrip yang perlu memeriksa apakah sebuah direktori
ada di tiga tempat berbeda. Tanpa fungsi, kamu menulis logika pemeriksaan
yang sama tiga kali. Dengan fungsi, kamu tulis sekali dan panggil tiga
kali. Kode lebih pendek, lebih mudah diperbaiki, dan lebih mudah dipahami.

Di Bash, fungsi didefinisikan tanpa daftar parameter eksplisit. Argumen
dikirim saat memanggil fungsi dan dibaca dengan \texttt{\$1},
\texttt{\$2}, dan seterusnya di dalam fungsi, sama seperti parameter
posisional di skrip. \texttt{\$\#} di dalam fungsi adalah jumlah argumen
fungsi, bukan argumen skrip.

Perintah \cmd{local} penting untuk mendeklarasikan variabel yang hanya
hidup di dalam fungsi. Tanpa \cmd{local}, variabel yang dibuat di dalam
fungsi mengubah variabel global dengan nama yang sama. Ini sumber bug
yang sulit dilacak.

Fungsi tidak bisa mengembalikan nilai seperti bahasa pemrograman lain.
Perintah \cmd{return} hanya mengembalikan angka 0 sampai 255 sebagai
status (seperti exit code). Untuk ``mengembalikan'' string atau angka,
cetak dengan \cmd{echo} di dalam fungsi dan tangkap hasilnya dengan
command substitution: \texttt{HASIL=\$(nama\_fungsi argumen)}.

File \textit{library} adalah file yang hanya berisi definisi fungsi,
tanpa kode yang dijalankan langsung. Skrip lain bisa memuat library ini
dengan perintah \cmd{source} atau \cmd{.} (titik). Pola ini mirip dengan
\cmd{import} di bahasa pemrograman lain: kamu mendefinisikan fungsi sekali
di satu tempat dan menggunakannya di banyak skrip.

\subsection*{Praktek 7.4: Refaktor Skrip Menggunakan Fungsi}

Tujuan praktek ini adalah memisahkan logika ke dalam fungsi dan
membuat file library yang bisa digunakan bersama.

\begin{enumerate}
    \item Buat file library dengan fungsi-fungsi utilitas umum:
\begin{longlisting}[language=bash, caption={Membuat file library.sh}]
nano ~/lab-os/chapter07-scripting/lib/library.sh
# Tulis isi berkas berikut
#!/bin/bash
# library.sh: Kumpulan fungsi utilitas untuk skrip sistem
# Cara pakai: source ~/lab-os/chapter07-scripting/lib/library.sh

# Tampilkan pesan log dengan timestamp
log_info() {
    echo "[$(date '+%H:%M:%S')] [INFO]  $*"
}

log_warn() {
    echo "[$(date '+%H:%M:%S')] [WARN]  $*" >&2
}

log_error() {
    echo "[$(date '+%H:%M:%S')] [ERROR] $*" >&2
}

# Keluar dengan pesan error
error_exit() {
    log_error "$1"
    exit "${2:-1}"
}

# Periksa apakah file ada dan bisa dibaca
file_bisa_dibaca() {
    [ -f "$1" ] && [ -r "$1" ]
}

# Periksa apakah direktori ada
dir_ada() {
    [ -d "$1" ]
}

# Konversi byte ke format mudah dibaca
bytes_ke_human() {
    local BYTES=$1
    if [ "$BYTES" -ge 1073741824 ]; then
        echo "$((BYTES / 1073741824)) GB"
    elif [ "$BYTES" -ge 1048576 ]; then
        echo "$((BYTES / 1048576)) MB"
    elif [ "$BYTES" -ge 1024 ]; then
        echo "$((BYTES / 1024)) KB"
    else
        echo "${BYTES} B"
    fi
}

# Pastikan direktori ada, buat jika belum
pastikan_dir() {
    if ! dir_ada "$1"; then
        mkdir -p "$1" && log_info "Direktori dibuat: $1"
    fi
}
\end{longlisting}

    \item Buat skrip utama yang menggunakan library:
\begin{longlisting}[language=bash, caption={Membuat skrip laporan-disk.sh yang menggunakan library}]
nano ~/lab-os/chapter07-scripting/skrip/laporan-disk.sh
# Tulis isi berkas berikut
#!/bin/bash
# laporan-disk.sh: Laporan penggunaan disk menggunakan library

LIB_PATH="$HOME/lab-os/chapter07-scripting/lib/library.sh"

if [ ! -f "$LIB_PATH" ]; then
    echo "Library tidak ditemukan: $LIB_PATH"
    exit 1
fi

source "$LIB_PATH"

OUTPUT_DIR="$HOME/lab-os/chapter07-scripting/logs"
pastikan_dir "$OUTPUT_DIR"

OUTPUT_FILE="$OUTPUT_DIR/laporan-disk-$(date '+%F').txt"

geneate_laporan() {
    log_info "Memulai laporan disk"
    echo "=== LAPORAN DISK ==="
    echo "Waktu: $(date '+%F %T')"
    echo "Host : $(hostname)"
    echo ""
    echo "--- Penggunaan Filesystem ---"
    df -h
    echo ""
    echo "--- Direktori Terbesar di /var ---"
    du -sh /var/* 2>/dev/null | sort -rh | head -5
    echo ""
    log_info "Laporan selesai: $OUTPUT_FILE"
}

geneate_laporan | tee "$OUTPUT_FILE"

chmod +x ~/lab-os/chapter07-scripting/skrip/laporan-disk.sh
./laporan-disk.sh
\end{longlisting}
    \textit{Amati:} skrip \cmd{laporan-disk.sh} tidak mendefinisikan
    sendiri fungsi seperti \cmd{log\_info} atau \cmd{pastikan\_dir}.
    Semua diambil dari library menggunakan \cmd{source}. Jika kamu perlu
    memperbaiki cara penulisan log, cukup ubah di \file{library.sh} dan
    semua skrip yang menggunakannya ikut berubah.

    \item Verifikasi bahwa fungsi library bisa dipanggil dari skrip lain:
\begin{lstlisting}[language=bash, caption={Menguji fungsi library secara langsung}]
source ~/lab-os/chapter07-scripting/lib/library.sh

log_info "Pengujian library berhasil"
log_warn "Ini pesan peringatan"
log_error "Ini pesan error (tampil di stderr)"

bytes_ke_human 1500
bytes_ke_human 2097152
bytes_ke_human 3221225472
\end{lstlisting}
\end{enumerate}

\begin{challengebox}
Pisahkan semua fungsi dari skrip \cmd{cek-disk.sh} di Praktek 7.3
ke dalam file library terpisah bernama \file{lib/library-disk.sh}. Buat
minimal tiga fungsi: \cmd{dapatkan\_persen\_disk()}, \cmd{status\_disk()}
yang mengembalikan teks ``PENUH''/``SEDANG''/``AMAN'' berdasarkan
persentase, dan \cmd{laporan\_semua\_disk()} yang mengiterasi semua
filesystem. Refaktor \cmd{cek-disk.sh} agar menggunakan ketiga fungsi
tersebut via \cmd{source}. Kemudian buat skrip utama baru bernama
\cmd{monitor-cepat.sh} yang meng-\cmd{source} kedua library
(\file{library.sh} dan \file{library-disk.sh}) dan menampilkan ringkasan
status disk dalam satu baris per filesystem.
\end{challengebox}

%% ------------------------------------------------------------------------
\section{Debugging dan Best Practices}
\label{sec:debugging-best-practices}

Skrip yang ditulis terburu-buru sering berjalan baik saat kondisi normal
tetapi gagal diam-diam saat ada sesuatu yang tidak terduga: file tidak
ada, variabel kosong, perintah gagal tanpa kamu sadari. Bash menyediakan
beberapa mekanisme untuk menangkap dan melaporkan kesalahan lebih awal.

Perintah \cmd{set -x} mengaktifkan mode \textit{trace}: setiap perintah
dicetak ke stderr sebelum dieksekusi, didahului dengan tanda \texttt{+}.
Ini sangat berguna untuk melihat alur eksekusi skrip yang tidak berjalan
sesuai harapan. Aktifkan dengan \texttt{bash -x skrip.sh} atau dengan
menambahkan \texttt{set -x} di dalam skrip.

Perintah \cmd{set -e} membuat skrip langsung berhenti saat ada perintah
yang gagal (exit code bukan 0). Tanpa \texttt{set -e}, Bash melanjutkan
eksekusi meskipun ada perintah yang gagal, yang sering menimbulkan
kerusakan lebih lanjut.

Perintah \cmd{set -u} membuat Bash melaporkan error jika kamu menggunakan
variabel yang belum didefinisikan. Tanpa \texttt{set -u}, variabel yang
belum didefinisikan diperlakukan sebagai string kosong, yang bisa
menimbulkan perilaku tak terduga.

Opsi \texttt{pipefail} membuat pipeline gagal jika salah satu komponen
di dalamnya gagal. Tanpa \texttt{pipefail}, exit code pipeline adalah
exit code perintah terakhir saja, sehingga kegagalan di tengah pipeline
tidak terdeteksi.

Kombinasi \texttt{set -euo pipefail} di awal skrip menangkap banyak
kelas kesalahan umum secara otomatis. Ini adalah standar untuk skrip
yang digunakan di produksi.

Perintah \cmd{trap} mendaftarkan perintah yang dijalankan saat sinyal
tertentu diterima atau saat skrip berakhir. Paling berguna adalah
\texttt{trap 'perintah-cleanup' EXIT}: perintah cleanup selalu dijalankan
saat skrip selesai, baik berhasil maupun gagal karena \texttt{set -e}.
Ini memastikan file sementara selalu dihapus dan tidak tertinggal.

ShellCheck adalah alat analisis statis yang memeriksa skrip Bash dan
melaporkan kesalahan umum, quoting yang salah, dan praktik yang tidak
dianjurkan. Pasang dengan \texttt{sudo apt install shellcheck} dan
jalankan dengan \texttt{shellcheck skrip.sh}.

\subsection*{Praktek 7.5: Debug Skrip Bermasalah}

Tujuan praktek ini adalah menggunakan teknik debugging untuk
menganalisis dan memperbaiki skrip, serta menambahkan mekanisme
keamanan ke skrip yang sudah ada.

\begin{enumerate}
    \item Buat skrip yang sengaja bermasalah untuk dianalisis:
\begin{lstlisting}[language=bash, caption={Membuat skrip bermasalah untuk latihan debug}]
nano ~/lab-os/chapter07-scripting/skrip/proses-log-bermasalah.sh
# Tulis isi berkas berikut
#!/bin/bash
# proses-log-bermasalah.sh: Skrip dengan bug tersembunyi

DIREKTORI=$1
FILE_TEMP=/tmp/proses-log-temp-$$

echo "Memproses log di: $DIREKTORI"
ls $DIREKTORI/*.log | wc -l
cp $DIREKTORI/*.log $FILE_TEMP
cat $FILE_TEMP
\end{lstlisting}

    \item Jalankan dengan \cmd{bash -x} untuk melihat setiap langkah
          eksekusi:
\begin{lstlisting}[language=bash, caption={Menggunakan mode trace untuk debugging}]
chmod +x ~/lab-os/chapter07-scripting/skrip/proses-log-bermasalah.sh

# Tanpa argumen (bug pertama)
bash -x ~/lab-os/chapter07-scripting/skrip/proses-log-bermasalah.sh

# Dengan direktori yang ada
bash -x ~/lab-os/chapter07-scripting/skrip/proses-log-bermasalah.sh \
  ~/lab-os/chapter07-scripting/logs
\end{lstlisting}
    \textit{Amati:} mode \texttt{-x} menampilkan setiap perintah sebelum
    dijalankan. Kamu bisa melihat nilai aktual variabel yang digunakan,
    bukan nama variabelnya.

    \item Buat versi yang sudah diperbaiki dengan \texttt{set -euo pipefail}
          dan \cmd{trap}:
\begin{longlisting}[language=bash, caption={Skrip yang sudah diperbaiki dengan mekanisme keamanan}]
nano ~/lab-os/chapter07-scripting/skrip/proses-log-aman.sh
# Tulis isi berkas berikut
#!/bin/bash
set -euo pipefail

# Trap untuk membersihkan file sementara saat skrip berakhir
FILE_TEMP="/tmp/proses-log-temp-$$"
cleanup() {
    if [ -f "$FILE_TEMP" ]; then
        rm -f "$FILE_TEMP"
        echo "[CLEANUP] File sementara dihapus: $FILE_TEMP"
    fi
}
trap cleanup EXIT

# Validasi argumen
if [ $# -ne 1 ]; then
    echo "Penggunaan: $0 <direktori-log>"
    exit 1
fi

DIREKTORI="$1"

# Validasi direktori
if [ ! -d "$DIREKTORI" ]; then
    echo "Error: direktori tidak ditemukan: $DIREKTORI"
    exit 2
fi

# Cari file log
JUMLAH_LOG=$(find "$DIREKTORI" -maxdepth 1 -name "*.log" | wc -l)

if [ "$JUMLAH_LOG" -eq 0 ]; then
    echo "Tidak ada file .log di: $DIREKTORI"
    exit 0
fi

echo "Ditemukan $JUMLAH_LOG file log di: $DIREKTORI"

# Gabungkan semua log ke file sementara
find "$DIREKTORI" -maxdepth 1 -name "*.log" -exec cat {} + > "$FILE_TEMP"

echo "Total baris dari semua log: $(wc -l < "$FILE_TEMP")"
echo "5 baris pertama:"
head -5 "$FILE_TEMP"

chmod +x ~/lab-os/chapter07-scripting/skrip/proses-log-aman.sh

# Buat beberapa file log untuk pengujian
echo -e "INFO: skrip mulai\nWARN: disk 75%\nINFO: selesai" \
    > ~/lab-os/chapter07-scripting/logs/test-a.log
echo -e "ERROR: koneksi gagal\nINFO: retry" \
    > ~/lab-os/chapter07-scripting/logs/test-b.log

./proses-log-aman.sh ~/lab-os/chapter07-scripting/logs
\end{longlisting}
    \textit{Amati:} meski skrip berhasil berjalan, pesan dari \cmd{cleanup}
    tetap muncul karena \cmd{trap ... EXIT} selalu dijalankan saat skrip
    berakhir, baik sukses maupun gagal.

    \item Uji skrip dengan berbagai kondisi error:
\begin{lstlisting}[language=bash, caption={Menguji penanganan berbagai kondisi error}]
# Tanpa argumen
./proses-log-aman.sh
echo "Exit code: $?"

# Direktori tidak ada
./proses-log-aman.sh /direktori-tidak-ada
echo "Exit code: $?"

# Direktori tanpa file log
mkdir -p /tmp/dir-kosong
./proses-log-aman.sh /tmp/dir-kosong
echo "Exit code: $?"
\end{lstlisting}
\end{enumerate}

\begin{challengebox}
Ambil skrip \cmd{cek-disk.sh} dari Praktek 7.3 dan tingkatkan dengan
tambahan berikut. Pertama, tambahkan \texttt{set -euo pipefail} di baris
kedua setelah shebang. Kedua, tambahkan \texttt{trap} yang membuat file
sementara \texttt{/tmp/cek-disk-\$\$.tmp} di awal skrip dan menghapusnya
saat exit. Ketiga, gunakan file sementara tersebut untuk menyimpan output
\cmd{df} sebelum diproses, sehingga skrip hanya memanggil \cmd{df} sekali.
Keempat, tambahkan opsi \texttt{-v} (verbose) menggunakan \cmd{getopts}
yang jika aktif menampilkan nama perangkat lengkap di samping setiap baris
status. Ini adalah pendekatan yang mirip dengan skrip \cmd{monitor-memori.sh}
yang akan kamu temukan di Bab~8.
\end{challengebox}

%% ------------------------------------------------------------------------
\section{Rangkuman}
\label{sec:rangkuman-bash-programming}

Dalam bab ini, kamu telah mempelajari:

\begin{itemize}[noitemsep,topsep=2pt]
    \item \textbf{Shebang dan izin eksekusi}: baris \texttt{\#!/bin/bash}
          menentukan interpreter. Perintah \texttt{chmod +x} memberikan
          izin eksekusi. Exit code 0 berarti sukses, bukan 0 berarti
          gagal.
    \item \textbf{Variabel Bash}: semua variabel adalah string. Tidak
          boleh ada spasi di sekitar \texttt{=}. Gunakan
          \texttt{"\$\{VARIABEL\}"} untuk mengakses nilai dengan aman.
    \item \textbf{Parameter posisional}: \texttt{\$1} sampai \texttt{\$9}
          adalah argumen skrip. \texttt{\$\#} jumlah argumen, \texttt{\$@}
          semua argumen sebagai daftar, \texttt{\$?} exit code terakhir.
    \item \textbf{Aritmetika}: gunakan \texttt{\$(( ekspresi ))} untuk
          perhitungan bilangan bulat. Gunakan \cmd{bc} untuk desimal.
    \item \textbf{Struktur kontrol}: \cmd{if}/\cmd{elif}/\cmd{else} untuk
          percabangan, \cmd{for} dan \cmd{while} untuk perulangan,
          \cmd{case} untuk mencocokkan pola nilai.
    \item \textbf{Fungsi}: didefinisikan dengan \texttt{nama() \{ ...\}},
          menerima argumen via \texttt{\$1}/\texttt{\$2}. Gunakan
          \cmd{local} untuk variabel lokal. Kembalikan nilai dengan
          \cmd{echo} dan tangkap dengan command substitution.
    \item \textbf{Library}: pisahkan fungsi ke file terpisah dan muat
          dengan \cmd{source}. Satu definisi, dipakai banyak skrip.
    \item \textbf{Debugging}: \texttt{bash -x} untuk trace, \texttt{set
          -euo pipefail} untuk menangkap kesalahan umum otomatis,
          \cmd{trap ... EXIT} untuk cleanup yang terjamin.
\end{itemize}

%% ------------------------------------------------------------------------
\section{Latihan}
\label{sec:latihan-bash-prog}

\textbf{Instruksi Umum}: kerjakan seluruh latihan di
\dir{~/lab-os/chapter07-scripting}. Simpan skrip di \dir{skrip}, library
di \dir{lib}, dan semua output di \dir{logs}. Setiap skrip harus
menyertakan shebang, komentar deskripsi, dan penanganan argumen yang
benar.

\subsubsection*{Latihan 7.1 Skrip Absensi dengan Laporan}
Buat skrip \cmd{absensi.sh} yang mencatat kehadiran dari command line.
Skrip menerima dua argumen: nama peserta dan status (\texttt{hadir},
\texttt{izin}, atau \texttt{alpha}). Setiap entri disimpan ke file
\texttt{logs/absensi-TANGGAL.txt} dengan format
\texttt{[HH:MM] NAMA - STATUS}. Tambahkan opsi \texttt{-r} yang
menampilkan rekapitulasi: jumlah hadir, izin, dan alpha dari file absensi
hari ini. Tambahkan opsi \texttt{-h} untuk menampilkan bantuan penggunaan.
Gunakan \cmd{getopts} untuk memproses opsi. Rekam minimal 6 entri dengan
status berbeda dan jalankan opsi \texttt{-r} untuk menampilkan
rekapitulasinya.


\subsubsection*{Latihan 7.2 Library dan Skrip Backup}
Buat file \file{lib/library-backup.sh} yang berisi minimal empat fungsi:
\cmd{buat\_nama\_arsip()} yang menghasilkan nama file dengan timestamp,
\cmd{validasi\_sumber()} yang memeriksa apakah direktori sumber ada,
\cmd{backup\_tar()} yang membuat arsip terkompresi, dan
\cmd{laporan\_backup()} yang menampilkan ukuran arsip dan jumlah file
di dalamnya. Buat skrip \cmd{backup-otomatis.sh} yang meng-\cmd{source}
library tersebut dan mendukung opsi \texttt{-s <sumber>}, \texttt{-t
<tujuan>}, \texttt{-v} (verbose), dan \texttt{-h} (bantuan) menggunakan
\cmd{getopts}. Skrip harus menggunakan \texttt{set -euo pipefail} dan
\cmd{trap} untuk cleanup file sementara. Uji dengan minimal dua skenario:
backup berhasil dan backup gagal karena sumber tidak ada.


\subsubsection*{Latihan 7.3 Skrip Monitor Sistem dengan Debug}
Buat skrip \cmd{monitor-sistem.sh} menggunakan template profesional
berikut: shebang, \texttt{set -euo pipefail}, \cmd{trap} cleanup,
\cmd{getopts} untuk opsi \texttt{-i <interval-detik>}, \texttt{-n
<jumlah-iterasi>}, dan \texttt{-o <file-output>}. Skrip memantau
penggunaan CPU, memori, dan disk setiap \texttt{interval} detik sebanyak
\texttt{n} kali, menyimpan setiap pembacaan ke file output dengan
timestamp. Gunakan fungsi terpisah untuk setiap komponen pemantauan.
Tambahkan mode debug: jika variabel lingkungan \var{DEBUG} bernilai
\texttt{1}, aktifkan \cmd{set -x} secara otomatis. Uji skrip dengan
\texttt{DEBUG=1 ./monitor-sistem.sh -i 2 -n 3} dan amati output trace
untuk memahami alur eksekusi.


%% ------------------------------------------------------------------------
\section{Referensi}
\label{sec:referensi-bash-programming}

\begin{itemize}
    \item Free Software Foundation. \textit{GNU Bash Reference Manual}.
          \url{https://www.gnu.org/software/bash/manual/}
    \item Shotts, William. \textit{The Linux Command Line}, 2nd Edition.
          No Starch Press, 2019.
    \item Robbins, Arnold dan Beebe, Nelson H.F. \textit{Classic Shell
          Scripting}. O'Reilly Media, 2005.
    \item Blum, Richard dan Bresnahan, Christine. \textit{Linux Command
          Line and Shell Scripting Bible}, 4th Edition. Wiley, 2021.
    \item ShellCheck: Alat analisis statis untuk skrip Bash.
          \url{https://www.shellcheck.net}
\end{itemize}
